\documentclass[11pt]{article}
\usepackage{mathblatt}
\begin{document}
\blattkopf*{Basisaufgaben}{BAS-Z1}{Basisaufgaben · Zettel · BAS-Z1}
\einheitenkopf[][Zettel 1]{Basisaufgaben · Zettel 1}
\zweigzeile{ohne Taschenrechner · je Aufgabe 1 Punkt · Lösungen auf der Rückseite}
% \small: zehn Aufgaben mit bis zu zwei Koordinatensystemen auf einer Seite (Render 28.09.)
\small

% 1: Bruchteil einer Fläche bestimmen – brueche-dezimalzahlen-basis-k1-v1 (Original 2026-FOR-B1b)
\begin{aufgabe}{Bei welcher Figur ist genau $\frac{1}{4}$ grau? \hfill (P26F) \\ \kreuz{P} \kreuz{Q} \kreuz{R} \kreuz{S}}

P \bruchkreis[0.8]{1}{3} \quad Q \bruchrechteck[2.5]{2}{8} \quad R \kreissektor[0.8]{120}{} \quad S \bruchrechteck[2.5]{2}{5}
\end{aufgabe}

% 2: Bruchteil einer Größe berechnen – brueche-dezimalzahlen-basis-k2-v1 (Original 2018-OS-B1a)
\begin{aufgabe}{Wie viel Gramm sind $\frac{2}{3}$ von $1{,}5$ kg? \hfill (P18) \\ \leerfeld[g]}
\end{aufgabe}

% 3: Zahlen in verschiedenen Darstellungen vergleichen – brueche-dezimalzahlen-basis-k4-v1 (Original 2023-OS-B1f)
\begin{aufgabe}{Welcher Wert ist der kleinste: $3{,}3$; $0{,}33$; $0{,}3^2$; $33\,\%$? \hfill (P23) \\ \leerfeld}
\end{aufgabe}

% 4: Termwert berechnen – rationale-zahlen-basis-k1-v1 (Original 2026-FOR-B1g)
\begin{aufgabe}{Berechne den Wert von $4 \cdot (x - 2)$ für $x = -3$. \hfill (P26F) \\ \leerfeld}
\end{aufgabe}

% 5: Exponent einer Potenz bestimmen – potenzen-wurzeln-basis-k1-v1 (Original 2024-OS-B1g)
\begin{aufgabe}{$2^x = 32$ – welche Zahl ist $x$? \hfill (P24) \\ x = \leerfeld}
\end{aufgabe}

% 6: Prozentwert berechnen – prozentrechnung-basis-k4-v1 (Original 2026-FOR-B1a)
\begin{aufgabe}{$20\,\%$ von $45\,€$ \hfill (P26F) \\ \leerfeld[€]}
\end{aufgabe}

% 7: Antiproportionale Zuordnung Dreisatz – zuordnungen-basis-k1-v1 (Original 2014-GYM-B1c)
\begin{aufgabe}{Ein Heuvorrat reicht für 5 Kühe 12 Tage. Wie viele Tage reicht er für 6 Kühe? \hfill (P14G) \\ \leerfeld[Tage]}
\end{aufgabe}

% 8: Lineare Gleichung lösen – lineare-gleichungen-basis-k1-v1 (Original 2024-OS-B1d)
\begin{aufgabe}{Löse die Gleichung $3 \cdot (x - 4{,}5) = 0$. \hfill (P24) \\ x = \leerfeld}
\end{aufgabe}

% 9: Lösung durch Einsetzen prüfen – quadratische-gleichungen-basis-k1-v1 (Original 2025-OS-B1h)
\begin{aufgabe}{Welcher Wert erfüllt $x \cdot (x + 6) = -8$? \hfill (P25) \\ \kreuz{$x = 2$} \kreuz{$x = 4$} \kreuz{$x = -2$} \kreuz{$x = -8$}}
\end{aufgabe}

% 10: Dreiecksungleichung anwenden – winkel-dreiecke-basis-k1-v1 (Original 2020-OS-B1i)
\begin{aufgabe}{Ein Dreieck ABC hat die Seite $a = 9$ cm. Was gilt für die Summe der beiden anderen Seiten $b$ und $c$? \hfill (P20) \\ \kreuz{$b + c > 9$ cm} \kreuz{$b + c = 9$ cm} \kreuz{$b + c < 9$ cm}}
\end{aufgabe}

\begleitteil
\erg{1}{Q, denn $\frac{2}{8} = \frac{1}{4}$}
\erg{2}{$1{,}5$ kg $= 1\,500$ g, $1\,500 : 3 \cdot 2 = 1\,000$ g}
\erg{3}{$0{,}3^2 = 0{,}09$; die anderen sind $3{,}3$ und zweimal $0{,}33$}
\erg{4}{$4 \cdot (-3 - 2) = 4 \cdot (-5) = -20$}
\erg{5}{$2, 4, 8, 16, 32$: fünf Faktoren, also $x = 5$}
\erg{6}{$0{,}2 \cdot 45 = 9\,€$}
\erg{7}{$5 \cdot 12 = 60$ Kuhtage; $60 : 6 = 10$ Tage}
\erg{8}{$x - 4{,}5 = 0$, also $x = 4{,}5$}
\erg{9}{$x = -2$, denn $(-2) \cdot (-2 + 6) = (-2) \cdot 4 = -8$}
\erg{10}{$b + c > 9$ cm}
\end{document}
